\section*{अध्याय \ref{ch:g7:triangles} --- त्रिभुज और कोण}

\begin{solution}{exo:g7:triangles:1}
पूरक: $60^\circ$; $45^\circ$; $18^\circ$।
संपूरक: $150^\circ$; $135^\circ$; $108^\circ$।
\end{solution}

\begin{solution}{exo:g7:triangles:2}
\omterm{def:g7:triangles:pairs}{शीर्षाभिमुख कोण} भी $35^\circ$ का है
(\cref{thm:g7:triangles:equalities}); बाक़ी दोनों कोण उसके संपूरक
हैं: $180 - 35 = 145^\circ$ प्रत्येक।
\end{solution}

\begin{solution}{exo:g7:triangles:3}
(a) $180 - 100 = 80^\circ$. \quad
(b) $180 - 118 = 62^\circ$. \quad
(c) $180 - 150 = 30^\circ$.
\end{solution}

\begin{solution}{exo:g7:triangles:4}
शिखर $40^\circ$: आधार वाले दोनों कोण $180 - 40 = 140^\circ$ को
बराबर बाँट लेते हैं: प्रत्येक $70^\circ$।

आधार कोण $70^\circ$: आधार वाले दोनों कोण बराबर हैं, इसलिए शिखर कोण
$180 - 2 \times 70 = 40^\circ$।
\end{solution}

\begin{solution}{exo:g7:triangles:5}
$(6, 8, 12)$: $12 < 6 + 8 = 14$: त्रिभुज बनता है।

$(5, 5, 11)$: $11 > 5 + 5 = 10$: असंभव।

$(4, 9, 13)$: $13 = 4 + 9$: चपटा ``त्रिभुज'' --- बिंदु एक ही रेखा
पर हैं, सच्चा त्रिभुज नहीं।

$(7, 7, 7)$: $7 < 14$: \omterm{def:g6:shapes:triangles}{समबाहु त्रिभुज}।
\end{solution}

\begin{solution}{exo:g7:triangles:6}
$x + 2x + 3x = 6x = 180^\circ$, इसलिए $x = 30^\circ$: कोण
$30^\circ$, $60^\circ$ और $90^\circ$ हैं --- एक \omterm{def:g6:shapes:triangles}{समकोण त्रिभुज}।
\end{solution}

\begin{solution}{exo:g7:triangles:7}
पहले से बताया गया माप: $\widehat{BAC} = 180 - (50 + 60) = 70^\circ$;
बनी हुई आकृति पर चाँदा इसकी पुष्टि कर देता है।
\end{solution}

\begin{solution}{exo:g7:triangles:8}
पहले कटान पर चारों कोण $54^\circ$, $126^\circ$, $54^\circ$,
$126^\circ$ हैं (शीर्षाभिमुख जोड़े और संपूरक)। \omterm{def:g6:lines:perp}{समांतर} रेखा दूसरे
कटान पर वही चारों माप दोहरा देती है (संगत कोण): कुल आठ कोण, चार
$54^\circ$ के और चार $126^\circ$ के।
\end{solution}

\begin{solution}{exo:g7:triangles:9}
मान लो तीसरी भुजा $c$ है। त्रिभुज असमिका माँगती है कि
$c < 8 + 3 = 11$ हो और साथ ही $8 < c + 3$, यानी $c > 5$। तो
$5 < c < 11$: $7$ cm की लंबाई चलती है; $12$ cm (या $4$ cm) की
नहीं चलती।
\end{solution}

\begin{solution}{exo:g7:triangles:10}
$\widehat A + \widehat B = 180 - 90 = 90^\circ$, और
$\widehat A = 2\widehat B$, इसलिए $3\widehat B = 90^\circ$:
$\widehat B = 30^\circ$ और $\widehat A = 60^\circ$।
\end{solution}

\begin{solution}{exo:g7:triangles:11}
विकर्ण $[AC]$, $ABCD$ को त्रिभुजों $ABC$ और $ACD$ में बाँट देता है।
\omterm{def:g4:geometry:polygon}{चतुर्भुज} के चारों कोण ठीक वही छह कोण हैं जो दोनों त्रिभुजों के हैं,
बस नए सिरे से जुड़े हुए ($A$ और $C$ वाले कोण दो-दो हिस्सों में बँट
जाते हैं)। कुल: $180 + 180 = 360^\circ$। \omterm{def:g4:geometry:polygon}{पंचभुज} एक \omterm{def:g5:solids:def}{शीर्ष} से तीन
त्रिभुजों में बँटता है: $3 \times 180 = 540^\circ$।
\end{solution}

\begin{solution}{pb:g7:triangles:1}
\textbf{1.} \omterm{def:g4:geometry:polygon}{षट्भुज} के एक \omterm{def:g5:solids:def}{शीर्ष} से खींचे गए विकर्ण उसे $4$
त्रिभुजों में काट देते हैं, इसलिए भीतरी कोणों का कुल
$4 \times 180 = 720^\circ$ है।

\textbf{2.} एक \omterm{def:g5:solids:def}{शीर्ष} से विकर्ण उस \omterm{def:g5:solids:def}{शीर्ष} और उसके दोनों पड़ोसियों
को छोड़कर बाक़ी सारे \omterm{def:g5:solids:def}{शीर्षों} तक जाते हैं: $n - 3$ विकर्ण, जो
\omterm{def:g4:geometry:polygon}{बहुभुज} को $n - 2$ त्रिभुजों में काटते हैं। \omterm{def:g4:geometry:polygon}{बहुभुज} का प्रत्येक
भीतरी कोण इन्हीं त्रिभुजों में बँट जाता है, इसलिए \omterm{def:g1:addition:def}{योगफल} है
\[
(n - 2) \times 180^\circ .
\]
जाँच: $n = 3$: $180^\circ$; $n = 4$: $360^\circ$
(\cref{exo:g7:triangles:11}); $n = 5$: $540^\circ$; $n = 6$:
$720^\circ$।

\textbf{3.} दशभुज: $8 \times 180 = 1\,440^\circ$। द्वादशभुज:
$10 \times 180 = 1\,800^\circ$।

\textbf{4.} $n$ से भाग देने पर: त्रिभुज $60^\circ$; वर्ग
$90^\circ$; \omterm{def:g4:geometry:polygon}{पंचभुज} $\frac{540}{5} = 108^\circ$; \omterm{def:g4:geometry:polygon}{षट्भुज}
$\frac{720}{6} = 120^\circ$; \omterm{def:g4:geometry:polygon}{अष्टभुज} $\frac{1080}{8} =
135^\circ$।

\textbf{5.} एक \omterm{def:g5:solids:def}{शीर्ष} बढ़ने से पंखे में एक त्रिभुज और जुड़ जाता है:
नए \omterm{def:g4:geometry:polygon}{बहुभुज} को $n - 1$ त्रिभुज चाहिए जहाँ पुराने को $n - 2$ चाहिए
थे। एक त्रिभुज ज़्यादा, यानी $180^\circ$ ज़्यादा।

\textbf{6.} चलने वाला संपूरक कोण से मुड़ता है:
$180^\circ - \widehat A$। $40$--$60$--$80$ वाले त्रिभुज में घुमाव
$140^\circ$, $120^\circ$ और $100^\circ$ हैं।

\textbf{7.} $140 + 120 + 100 = 360^\circ$: एक पूरा चक्कर।

\textbf{8.} पूरे चक्कर में तुम्हारी नाक जिस दिशा से शुरू होती है
उसी दिशा में लौट आती है, और बीच में ठीक एक बार चारों ओर घूम जाती
है: कोनों पर किया गया सारा मुड़ना मिलकर एक पूरा चक्कर,
$360^\circ$, बनता है --- कोनों की संख्या और मुहल्ले की शक्ल चाहे
जो हो। (यह बात $n$ से बिलकुल आज़ाद है, और यही इसकी ख़ूबसूरती है।)

\textbf{9.} प्रत्येक कोने पर भीतरी कोण $+$ घुमाव $= 180^\circ$।
$n$ कोनों पर जोड़ने से:
\[
\text{(भीतरी कोणों का योग)} + 360^\circ = 180^\circ \times n,
\qquad\text{इसलिए}\qquad
\text{भीतरी कोणों का योग} = (n - 2) \times 180^\circ
\]
--- सवाल 2 वाला सूत्र, चलकर दोबारा सिद्ध हुआ।

\textbf{10.} भीतरी कोण $150^\circ$ का मतलब है कि प्रत्येक घुमाव
$30^\circ$ है, और $n = \frac{360}{30} = 12$: सम द्वादशभुज।
भीतरी कोण $155^\circ$ का मतलब होता $25^\circ$ के घुमाव, और
$\frac{360}{25} = 14.4$ कोनों की पूरी संख्या नहीं है: ऐसा कोई सम
\omterm{def:g4:geometry:polygon}{बहुभुज} है ही नहीं।

\textbf{11.} त्रिभुज ($60^\circ$): $6$ मिलते हैं,
$6 \times 60 = 360$। वर्ग ($90^\circ$): $4$ मिलते हैं। \omterm{def:g4:geometry:polygon}{षट्भुज}
($120^\circ$): $3$ मिलते हैं, $3 \times 120 = 360$। तीनों पूरी
तरह बंद हो जाते हैं।

\textbf{12.} \omterm{def:g4:geometry:polygon}{पंचभुज} के तीन कोने: $3 \times 108 = 324^\circ$ ---
$36^\circ$ की ख़ाली जगह बच जाती है। चार कोने:
$4 \times 108 = 432^\circ > 360^\circ$ --- चढ़ाव। दोनों में से
कुछ नहीं चलता: सम \omterm{def:g4:geometry:polygon}{पंचभुज} समतल नहीं ढक सकते।

\textbf{13.} ढकाई के किसी कोने पर कम से कम $3$ टाइलें मिलती हैं।
छह से ज़्यादा भुजाओं वाले सम \omterm{def:g4:geometry:polygon}{बहुभुज} के भीतरी कोण $120^\circ$ से
\emph{बड़े} होते हैं (सवाल 4, और $n$ के साथ कोण का बढ़ना), इसलिए
तीन कोने मिलकर $3 \times 120 = 360^\circ$ से ऊपर निकल जाते हैं:
चढ़ाव, यानी असंभव। ठीक छह भुजाओं पर $3 \times 120 = 360$ चलता है;
छह से नीचे के मामले सवाल 11 और 12 निपटा देते हैं। सम ढकने वालों
की पूरी सूची: त्रिभुज, वर्ग, \omterm{def:g4:geometry:polygon}{षट्भुज}।

\textbf{14.} \omterm{def:g4:geometry:polygon}{अष्टभुज} और वर्ग: $135 + 135 + 90 = 360^\circ$: कोना
बंद हो जाता है --- गलियों वाला जाना-पहचाना नमूना। वर्ग, \omterm{def:g4:geometry:polygon}{षट्भुज},
द्वादशभुज: द्वादशभुज का घुमाव $\frac{360}{12} = 30^\circ$ है,
इसलिए उसका भीतरी कोण $150^\circ$, और
$90 + 120 + 150 = 360^\circ$: यह भी चलता है।

\textbf{15.} छत्ते के कोने: $3 \times 120 = 360^\circ$, बिलकुल
ठीक। बराबर \omterm{def:g6:measure:area}{क्षेत्रफल} वाली त्रिभुज, वर्ग और \omterm{def:g4:geometry:polygon}{षट्भुज} टाइलों में
\omterm{def:g4:geometry:polygon}{षट्भुज} की सीमा सबसे छोटी होती है --- तीनों में वही रानी दीदो के
वृत्त के सबसे पास है (\cref{pb:g6:measure:1}) --- इसलिए षट्भुजाकार
ख़ाने उतने ही शहद को सबसे कम मोम में बंद कर लेते हैं: मधुमक्खियाँ
ज्यामिति जानने वालों की तरह फ़र्श ढकती हैं।

\end{solution}
